\section{Adaptive comparison and an unlimited supply of slots}
\label{sec:comparison}

Our goal is to compare a cut of the defect of a central matrix cone
with one of its diagonal corners, with a witness bound independent of
matrix size.  Norm visibility supplies finite coordinate ideal
multiplicities $L_n$, which need not be bounded.  We fix the norm
maps and these integers first, and then choose a tracial cone with
relative defect error $\eta_n/L_n$, where $\eta_n\to0$.
The ideal comparison bounds the controlling rank by a multiple of
$L_n$ times the target rank; passing to a diagonal corner adds only
the fixed matrix-size factor.  The reciprocal $L_n$ in the trace
error cancels this multiplicity.  We must therefore obtain an exact
coordinate Cuntz inequality before using it.  A further spectral
margin bounds the witnesses independently of matrix size, allowing
the final construction of countably many orthogonal equivalent
central slots.

Fix a countable norm-dense family $\mathcal A$ in the positive
contractions of $A$, closed under positive rational spectral cuts.
Such a family is obtained by starting with a countable dense family
and adjoining its iterated rational cuts.  Enumerate the tests
\[
 (a_r,\eps_r)\in\mathcal A\times(\mathbb Q\cap(0,1)),
 \qquad r\ge1.
\]
We shall use the fixed scalar functions
\begin{equation}\label{cmp:cutoff-functions}
 f(t)=\min(1,2t),\qquad g(t)=\max(0,2t-1)
       \quad(0\le t\le1).
\end{equation}
Thus $g$ is nondecreasing, vanishes on $[0,1/2]$, and $g(1)=1$;
also $f(t)\ge t$ and $fg=g$.  Functional calculus on an order-zero
map has the meaning fixed in the preliminaries.  The function $f$
enlarges its weight, while $g$ selects a part on which $f$ is one;
the enlarged cone will consequently fix the slots cut out by $g$.

\begin{proposition}[Coordinate ideal visibility]\label{prop:idealvisibility}
For every fixed $k\ge2$ there is a \cpc order-zero map
$\nu:M_k\to A_\omega\cap A'$, represented by \cpc maps
$\nu_n:M_k\to A$, with weights $h_n^\nu=\nu_n(1)$, such that the
following holds.  For every test $(a,\eps)$, put
\[
 b=(a-\eps/8)_+,\qquad
 y_n^\nu=a^{1/2}g(h_n^\nu)^2a^{1/2}.
\]
There are integers $L_{a,\eps,n}\ge1$ and an $\omega$-large set
$E_{a,\eps}$ on which
\begin{equation}\label{cmp:I}
 [b]\le L_{a,\eps,n}[(y_n^\nu-\eps/20)_+]
          \quad\text{in }\Cu(A).
\end{equation}
The integers need not be bounded as $n$ varies.
\end{proposition}

\begin{proof}
First consider finitely many tests.  For one of them apply
Proposition~\ref{prop:norm} with
\[
 x=(a-\eps/16)_+,
 \qquad \gamma=\eps/64.
\]
It gives a norm-central order-zero map with positive contractive
weight representatives $t_n$ and a constant $0<\kappa\le1$ such
that, on a common $\omega$-large set,
\begin{equation}\label{cmp:norm-trigger}
 \|\sigma(x)\|\ge\gamma
 \quad\Longrightarrow\quad
 \|\sigma(xt_n)\|\ge\kappa
 \qquad\text{for every irreducible }\sigma.
\end{equation}
If no representation satisfies the trigger, one can take the zero
map and any $0<\kappa\le1$; the assertions below then remain
vacuous on the relevant representations.  Choose a continuous
positive contraction $\chi$ on $[0,1]$, zero at zero and equal to
one on $[\kappa/2,1]$, and reparametrize the order-zero map by $\chi$.

Do this for each of the finitely many tests, choosing commuting
maps by successive reindexing.  Here reindexing preserves the
representation quantifier in \eqref{cmp:norm-trigger}.  More
explicitly, at coordinate $n$, previously chosen matrix-unit lifts
form a finite list of elements of $A$.  Choose an index of the next
map at which its norm commutators with those elements and the first
$n$ dense tests are small, its order-zero defects are small, and
its own uniform implication \eqref{cmp:norm-trigger} holds.
Each condition holds on an $\omega$-large set, and a finite
intersection is nonempty.  The implication is simultaneous over
all representations at the selected index, so it is retained
without choosing representations during reindexing.  Join the
reparametrized maps using Lemma~\ref{lem:join}.

Let $\nu$ denote this finite join for the moment, and choose \cpc
coordinate lifts with weights $h_n^\nu$.  We claim that, on an
$\omega$-large set, simultaneously in all irreducible representations,
\begin{equation}\label{cmp:visibility-implication}
 \|\sigma(a)\|>3\eps/32
 \quad\Longrightarrow\quad
 \|\sigma(a^{1/2}g(h_n^\nu)^2a^{1/2})\|>\eps/20.
\end{equation}
Suppose instead that its failure set belongs to $\omega$, and
choose a failing representation $\sigma_n$ at each such coordinate.
Intersect with the good set for \eqref{cmp:norm-trigger}.  Since
\[
 \|\sigma_n((a-\eps/16)_+)\|
   >3\eps/32-\eps/16=\eps/32>\gamma,
\]
that trigger applies, even if the represented norms tend to the
boundary $3\eps/32$.

Take the Hilbert space ultraproduct of these representations, and
write $\bar a,\bar t,\bar H$ for the resulting positive operators
representing $a$, the old unmodified weight for this test, and the
joined weight.  Norm centrality gives $[\bar a,\bar t]=0$, and
\eqref{cmp:norm-trigger} gives
\[
 \| (\bar a-\eps/16)_+\bar t\|\ge\kappa.
\]
Consequently the joint spectral part where $\bar a>\eps/16$ and
$\bar t>\kappa/2$ is nonzero.  Indeed the product is zero where
$\bar a\le\eps/16$ and has norm at most $\kappa/2$ where
$\bar t\le\kappa/2$.  On the indicated nonzero part the
reparametrized weight $\chi(\bar t)$ is one.  The joined weight
dominates this weight and commutes with it and with $\bar a$;
being a contraction, it too is one there.  It follows that
\[
 \|\bar a^{1/2}g(\bar H)^2\bar a^{1/2}\|
       \ge\eps/16>\eps/20.
\]
This contradicts the upper bound at the failing coordinates.
Arbitrary positive weight lifts of the joined map give the same
limiting operators, so no particular lifting convention was needed
for this contradiction.  This proves
\eqref{cmp:visibility-implication} for each test in the finite group.

For each $l$ perform this construction for the first $l$ tests.
Take \cpc coordinate lifts of the resulting maps.  Choose a decreasing
sequence of $\omega$-large sets on which the first $l$ norm
centrality tests and matrix-unit order-zero identities hold to
accuracy tending to zero, and all the first $l$ implications
\eqref{cmp:visibility-implication} hold exactly.  Include $\{n\ge l\}$
in the $l$th set, and use its largest admissible stage at coordinate
$n$.  Lemma~\ref{lem:diagonal} applies, since all these maps and
weights are contractions.  The diagonal defines a single
norm-central order-zero $\nu$, with \cpc lifts, and preserves each
of the implications on an $\omega$-large set.

At a coordinate where \eqref{cmp:visibility-implication} holds,
every irreducible representation annihilating
$(y_n^\nu-\eps/20)_+$ annihilates $(a-3\eps/32)_+$.  Indeed
annihilation of the former is precisely the condition
$\|\sigma(y_n^\nu)\|\le\eps/20$, which implies
$\|\sigma(a)\|\le3\eps/32$.  Primitive ideal separation gives
\[
 (a-3\eps/32)_+\in\Ideal_A((y_n^\nu-\eps/20)_+).
\]
The additional positive cut $\eps/8-3\eps/32=\eps/32$ now gives a
finite integer $L_{a,\eps,n}$ in \eqref{cmp:I}, by the compact-cut
ideal comparison from the preliminaries.  Choose the integer to be
at least one and put it equal to one outside the good set.  This
makes all constants specified at all coordinates and proves the
proposition.
\end{proof}

\begin{proposition}[An adaptive rank gap]\label{prop:rankgap}
For every fixed $k\ge2$ there is a \cpc order-zero map
$\zeta^k:M_k\to A_\omega\cap A'$ with weight $H$ such that, for every
test $(a,\eps)$, if
\begin{equation}\label{cmp:u-y}
 u^k(a)=(1-f(H))a,
 \qquad y^k(a)=g(\zeta^k)(e_{11})^2a,
\end{equation}
there is $r_{a,\eps}^k\in A_\omega$ satisfying
\begin{equation}\label{cmp:bounded-comparison}
 (u^k(a)-\eps)_+
   =(r_{a,\eps}^k)^*y^k(a)r_{a,\eps}^k,
 \qquad \|r_{a,\eps}^k\|^2\le32/\eps.
\end{equation}
In particular the witness bound is independent of $k$.
\end{proposition}

\begin{proof}
Apply Proposition~\ref{prop:idealvisibility} for this $k$.
\emph{Fix its coordinate maps $\nu_n$, their weights, and every
integer $L_{a,\eps,n}$ before making any subsequent choices.}
They will not be reindexed in this proof.

For a test $(a,\eps)$ retain $b=(a-\eps/8)_+$ and put
$d=(a-\eps/4)_+$.  Then $d\in\Ideal_A(b)$ and
$[d]\le[b]$, so Proposition~\ref{prop:localized-trace} applies to
$(d,b)$ with relative rank constant one.  Let $E_r$ be the good
sets in \eqref{cmp:I} for the enumerated tests.  Form decreasing
$\omega$-large sets
\[
 G_l=\{n\ge l\}\cap\bigcap_{r\le l}E_r,
\]
and let $l(n)$ be the largest $l\le n$ with $n\in G_l$, or zero
if there is none.  Thus $l(n)\to_\omega\infty$, and all ideal
bounds for the first $l(n)$ tests are valid at coordinate $n$.
Choose positive numbers $\eta_n\to0$, for example $\eta_n=1/n$.

At coordinate $n$ use the finite-coordinate conclusion of
Proposition~\ref{prop:localized-trace} for the first $l(n)$ pairs
$(d_r,b_r)$, with the individually prescribed accuracies
\[
 \eta_n/L_{a_r,\eps_r,n}.
\]
Include the first $n$ members of a fixed dense family in $A$, and
all matrix-unit values $\nu_n(e_{ij})$, in the commutator tests.
Ask for norm commutators and norm order-zero defects less than
$1/n$.  Denote the resulting \cpc map by $\theta_n$, and its weight
by $h_n^\theta$.  Where $l(n)=0$ one may use the zero map.  Every
finite problem here has already specified its finite integers;
its requested trace accuracy is allowed to be arbitrarily small.
Consequently the maps represent a norm-central order-zero map
$\theta$ commuting with $\nu$, and, for each fixed test, on an
$\omega$-large set,
\begin{equation}\label{cmp:adaptive-trace}
 \tau\bigl(d^{1/2}(1-h_n^\theta)d^{1/2}\bigr)
 \le\frac{\eta_n}{L_{a,\eps,n}}d_\tau(b)
       \qquad(d_\tau(b)<\infty).
\end{equation}
At zero rank this asserts exact zero; alternatively it follows
directly from $[d]\le[b]$.

Join $\theta$ and $\nu$ to obtain $\zeta^k$.  Use the following
specific positive contractive representatives of its weight:
\begin{equation}\label{cmp:positive-weight-lift}
 H_n=1-(1-h_n^\theta)^{1/2}(1-h_n^\nu)
                         (1-h_n^\theta)^{1/2}.
\end{equation}
They belong to $A$: equivalently, they are
$h_n^\theta+(1-h_n^\theta)^{1/2}h_n^\nu(1-h_n^\theta)^{1/2}$.
They lie between $h_n^\theta$ and $1$.  Since the two weights
commute in the ultrapower, they represent the joined weight $H$.
For $u^k(a)$ use the positive contractive representatives
\begin{equation}\label{cmp:source-lift}
 u_n=(1-f(H_n))^{1/2}a(1-f(H_n))^{1/2}.
\end{equation}
For $y^k(a)$ choose positive contractive representatives $y_n$;
for example, if $B_n$ are positive contractive representatives of
$g(\zeta^k)(e_{11})$, one can take $y_n=B_naB_n$.
All products in \eqref{cmp:u-y} are positive because the matrix
map is norm-central.

We first convert the norm visibility into an exact coordinate
Cuntz inequality.  In $A_\omega$ the old weight $h^\nu$ and $H$
commute, $h^\nu\le H$, and both commute with $a$.  Monotonicity of
$g$ in this commutative functional calculus gives
\begin{equation}\label{cmp:total-weight-domination}
 y^\nu=ag(h^\nu)^2\le ag(H)^2.
\end{equation}
Write $\varrho$ for the supporting homomorphism of $\zeta^k$.
Since $a$ commutes with the range of $\zeta^k$, it also commutes
with the supporting matrix units: these are the strong limits of
$\zeta^k(e_{ij})(H+t)^{-1}$ as $t\downarrow0$ in the bidual.
The right side is the orthogonal sum of
\[
 z_i=ag(H)^2\varrho(e_{ii})\quad(1\le i\le k),
 \qquad z_1=y^k(a).
\]
The products are in $A_\omega$.  All $z_i$ are Cuntz equivalent,
as are their corresponding positive cuts.  Indeed
$a^{1/2}g(H)\varrho(e_{ij})\in A_\omega$ has its two squared
moduli equal to $z_j,z_i$; taking polar cuts gives the same fact
for $(z_j-t)_+$ and $(z_i-t)_+$.  The two positives in
\eqref{cmp:total-weight-domination} commute, so scalar functional
calculus, splitting a cut of an orthogonal sum, and enlarging its
individual cuts give
\begin{equation}\label{cmp:ultrapower-cut}
 [(y^\nu-\eps/22)_+]
       \le k[(y^k(a)-\eps/30)_+]
       \quad\text{in }\Cu(A_\omega).
\end{equation}

Here the cut margins yield an exact inequality before any
multiplication by a varying integer.  Put
$\Delta=\eps/20-\eps/22>0$.
Choose one rectangular comparison witness for
\eqref{cmp:ultrapower-cut} in $M_{k,1}(A_\omega)$ with norm error
less than $\Delta/2$.  Such rectangular witnesses suffice because
the source is in one corner and the target is the diagonal of
$k$ elements of $A_\omega$.  Choose bounded coordinate lifts $X_n$
of this witness.  Since functional calculus commutes with the
quotient, on an $\omega$-large set,
\[
 \left\|(y_n^\nu-\eps/22)_+
       -X_n^*\operatorname{diag}\bigl((y_n-\eps/30)_+,
                           \ldots,(y_n-\eps/30)_+\bigr)X_n
 \right\|<\Delta.
\]
Norm comparison at that coordinate gives
\[
 [(y_n^\nu-\eps/20)_+]
 \le k[(y_n-\eps/30)_+]
 \le k[(y_n-\eps/32)_+].
\]
Intersect with the good set for \eqref{cmp:I}.  Only now multiply
the exact Cuntz inequality by $L_n=L_{a,\eps,n}$ to obtain
\begin{equation}\label{cmp:exact-coordinate-bound}
 [b]\le kL_n[(y_n-\eps/32)_+].
\end{equation}
The bound on $X_n$ may depend on the fixed $k$ and the test.  No
norm error is multiplied by $L_n$ in this argument.

We next prove the coordinate rank gap
\begin{equation}\label{cmp:R}
 d_\tau((u_n-\eps)_+)
       \le\tfrac12 d_\tau((y_n-\eps/32)_+)
       \qquad\text{for every }\tau
\end{equation}
on an $\omega$-large set.  Put
$c_n=(1-f(H_n))^{1/2}$ and $v_n=c_ndc_n$.
Since $\|a-d\|\le\eps/4$, one has
$\|u_n-v_n\|\le\eps/4$.  Comparing $u_n$ with
$(v_n-\eps/2)_+$ gives
\[
 [(u_n-\eps)_+]\le[(v_n-\eps/2)_+],
 \qquad
 d_\tau((u_n-\eps)_+)\le(2/\eps)\tau(v_n).
\]
For the first inequality the norm error is at most
$\eps/4+\eps/2<\eps$; the second is the elementary spectral
trace bound.  From the scalar inequality $f(t)\ge t$ and
\eqref{cmp:positive-weight-lift},
\[
 f(H_n)\ge H_n\ge h_n^\theta.
\]
This uses functional calculus on $H_n$ itself and positive
domination of $H_n$; it makes no operator-monotonicity assertion
about $f$.  Traciality and \eqref{cmp:adaptive-trace} consequently
give, whenever $d_\tau(b)<\infty$,
\begin{equation}\label{cmp:source-rank-estimate}
 d_\tau((u_n-\eps)_+)
 \le\frac{2}{\eps}
     \tau\bigl(d^{1/2}(1-f(H_n))d^{1/2}\bigr)
 \le\frac{2\eta_n}{\eps L_n}d_\tau(b).
\end{equation}

This estimate handles every possible trace value.  Set
$R_n=d_\tau((y_n-\eps/32)_+)$.  If $R_n<\infty$, the exact
inequality \eqref{cmp:exact-coordinate-bound} gives
$d_\tau(b)\le kL_nR_n<\infty$, and therefore
\[
 d_\tau((u_n-\eps)_+)
      \le (2k\eta_n/\eps)R_n.
\]
For the fixed $k$ and test, eventually $2k\eta_n/\eps\le1/2$.
If $R_n=0$, the same argument first forces $d_\tau(b)=0$ and then
gives exactly zero on the left.  If $R_n=\infty$, inequality
\eqref{cmp:R} is automatic in $[0,\infty]$.  All good sets used
in these alternatives were chosen uniformly over the trace
families; thus their intersection proves \eqref{cmp:R}
simultaneously for every lower semicontinuous extended trace.

Apply Lemma~\ref{lem:G} in $A$ at each coordinate of this common
good set.  It gives
\[
 (u_n-\eps)_+\precsim (y_n-s)_+,
 \qquad s=\eps/32.
\]
Choose $x_n\in A$ with
\[
 \|x_n^*(y_n-s)_+x_n-(u_n-\eps)_+\|<1/n.
\]
These witnesses need not themselves be bounded.  Define the
continuous scalar function
\[
 q_s(t)=
 \begin{cases}\bigl((t-s)_+/t\bigr)^{1/2},&t>0,\\0,&t=0,
 \end{cases}
\]
and set $r_n=q_s(y_n)x_n$ on the good set and $r_n=0$ elsewhere.
Then
\[
 r_n^*y_nr_n=x_n^*(y_n-s)_+x_n,
 \qquad
 \|r_n\|^2\le s^{-1}\|x_n^*(y_n-s)_+x_n\|
                 \le(32/\eps)(1+1/n).
\]
The norm estimate follows from $q_s(t)^2\le s^{-1}(t-s)_+$.
Thus these representatives have the fixed bound
$\|r_n\|^2\le64/\eps$, and their ultrapower element has norm
squared at most $32/\eps$.  The compression errors tend to zero
along $\omega$, proving \eqref{cmp:bounded-comparison}.  Neither
bound depends on $k$ or on any $L_n$.
\end{proof}

\begin{proposition}[Infinitely many central slots]\label{prop:slots}
For every $p\ge2$ there are a \cpc order-zero map
$\alpha:M_p\to A_\omega\cap A'$, with weight $h=\alpha(1)$,
and contractions $s_j\in A_\omega\cap A'$ for $j\ge1$ such that,
putting $P=s_1$, one has
\begin{equation}\label{cmp:slot-relations}
 P\ge0,\qquad s_i^*s_j=\delta_{ij}P^2,
 \qquad \alpha(e_{11})s_j=s_j.
\end{equation}
For every test $(a,\eps)$ there is $r_{a,\eps}\in A_\omega$ with
\begin{equation}\label{cmp:slot-comparison}
 ((1-h)a-\eps)_+
       =r_{a,\eps}^*(P^2a)r_{a,\eps},
 \qquad \|r_{a,\eps}\|^2\le64/\eps.
\end{equation}
No centrality is asserted or needed for these comparison witnesses.
\end{proposition}

\begin{proof}
For fixed $m\ge1$, apply Proposition~\ref{prop:rankgap} with $k=pm$
and identify $M_{pm}=M_p\otimes M_m$.  Define
\[
 \alpha_m(z)=f(\zeta^{pm})(z\otimes1_m),\qquad
 s_j^{(m)}=g(\zeta^{pm})(e_{11}^{(p)}\otimes e_{j1}^{(m)})
       \quad(1\le j\le m),
 \qquad P_m=s_1^{(m)}.
\]
These are norm-central contractions, and $P_m$ is positive.
The matrix-unit relations for the supporting homomorphism give
\[
 (s_i^{(m)})^*s_j^{(m)}=\delta_{ij}P_m^2.
\]
The scalar identity $fg=g$ gives
$\alpha_m(e_{11})s_j^{(m)}=s_j^{(m)}$.
The weight of $\alpha_m$ is $f(H^{pm})$.  Under the indicated
identification the first diagonal slot of $M_{pm}$ is
$e_{11}^{(p)}\otimes e_{11}^{(m)}$, so
Proposition~\ref{prop:rankgap} supplies, for every test,
\[
 ((1-\alpha_m(1))a-\eps)_+
       =(r_{a,\eps}^{(m)})^*(P_m^2a)r_{a,\eps}^{(m)}.
\]
The witness norm bounds are independent of $m$.

We spell out the diagonalization because the number of slots grows.
For every $m$ choose \cpc coordinate lifts $\alpha_{m,n}$, contractive
lifts $s_{j,m,n}$ for $1\le j\le m$, and positive contractive lifts
$P_{m,n}=s_{1,m,n}$.  Choose lifts of each comparison witness bounded
by $\sqrt{64/\eps}$; the sharper ultrapower bound in
Proposition~\ref{prop:rankgap} leaves room for this choice.  Put
$h_{m,n}=\alpha_{m,n}(1)$.  At stage $l$ use $m=l$ and require:
\begin{enumerate}
 \item the matrix-unit order-zero identities of $\alpha_{l,n}$,
       and its commutators with the first $l$ dense tests of $A$,
       to have error less than $1/l$;
 \item the centrality and the two algebraic relations in
       \eqref{cmp:slot-relations} for all slot indices at most $l$
       to have error less than $1/l$;
 \item for the first $l$ comparison tests, the equations with
       positive coordinate expressions
       \[
       \bigl((1-h_{l,n})^{1/2}a(1-h_{l,n})^{1/2}-\eps\bigr)_+
       \quad\text{and}\quad
       r_{a,\eps,l,n}^*P_{l,n}aP_{l,n}r_{a,\eps,l,n}
       \]
       to have error less than $1/l$.
\end{enumerate}
For fixed $l$ all these finitely many estimates hold on an
$\omega$-large set, because their equations are exact in the
ultrapower.  The positive coordinate expressions represent the
desired products since the weights and slots commute with $A$ in
the ultrapower.

Intersect the first $l$ such good sets with $\{n\ge l\}$ to obtain
decreasing $E_l\in\omega$, and take $m(n)$ to be the largest
$l\le n$ with $n\in E_l$ (zero if none).  At coordinate $n$ use
the row indexed by $m(n)$.  For a slot or witness whose index is
larger than $m(n)$, use zero at that coordinate.  Every fixed
variable is eventually present along $\omega$, since
$m(n)\to_\omega\infty$.  All maps and slots are contractions,
and the witness for a fixed test has its individual bound
$\sqrt{64/\eps}$ independently of the row.  These are precisely
the boundedness hypotheses of Lemma~\ref{lem:diagonal}.
The selected maps therefore define a \cpc norm-central order-zero
$\alpha$, and the selected slots and witnesses define elements
with \eqref{cmp:slot-relations}--\eqref{cmp:slot-comparison}.
No uniform trace estimate as $m$ grows was needed: each fixed
matrix size had already produced exact equations with the stated
individual norm bounds.
\end{proof}
